% proofs/08_blackwell.tex
% The tagged experiment under a tighter disclosure rule.
% Fragment, to be \input by appendix.tex.

\section{Information is monotone in the rule}
\label{sec:blackwell}

Let $\theta=(j(s),s)$ and fix a rule $r=(\tau,T)$. At date $d$, set
$D_r^d=\ind\{a=1,\ c(\theta;\tau)+T\leq d\}$ and $L_{r,e}=D_r^e$. At the control date the market
observes the tagged output
\[
Y_r=
\begin{cases}
(F,\mathsf S_{F,r}),&D_r^H=1,\\
(P,\mathcal H^P_r),&D_r^H=0,
\end{cases}
\qquad
\mathsf S_{F,r}=(B_r^F,Q_r^F,a{=}1),
\]
where
\[
\mathcal H^P_r=(X_0,\ldots,X_H;L_{r,0},\ldots,L_{r,H}).
\]
Write $\mathcal C_{F,r}=\{D_r^H=1\}$, $\mathcal C_{P,r}=\{D_r^H=0\}$, and
$\Omega_r=\Prb(D_r^H=1)$, used below in the strictness criterion.
Every flag coordinate in a pooled history is zero. Prices may be appended to either branch. They
are pinned measurable functions of the displayed information and do not enlarge the experiment.
Write $\mathcal E_r(\kappa)$ for the conditional channel from $\theta$ to $Y_r$ at liquidity
$\kappa$.

\csname unlabelledtrue\endcsname
\begin{theorem}[Tightening either disclosure dial]
\label{thm:blackwell}
Fix a common $\kappa\in[0,1]$ and fixed policies. Assume:
\begin{enumerate}[label=(B\arabic*),leftmargin=3.0em,itemsep=3pt]
\item \emph{Standing conditions.} Conditions \ref{p:space}--\ref{p:decoder} hold. Condition (S14),
  identified flagged types, is load-bearing: on the tighter rule's
  flagged region the tuple recovers the signal, and hence the type.
\item \emph{A tighter rule.} Either $b_0<\tau'<\tau$ at a common $T$, with
  $r_+=(\tau',T)$ and $r_-=(\tau,T)$, or $T'<T$ at a common $\tau$, with
  $r_+=(\tau,T')$ and $r_-=(\tau,T)$.
\end{enumerate}
Then the tighter tagged experiment is weakly more informative in Blackwell's order about
$\theta$ and about the price-relevant state $(v,a)$. Specifically, there is a Markov kernel $K_H$,
independent of the true type, such that
\begin{equation}
\label{eq:blackwell-kernel}
\mathcal L(Y_{r_-}\mid\theta)
=\int K_H(y,\mathord\cdot)\,\mathcal L(Y_{r_+}\in\mathrm dy\mid\theta).
\end{equation}
Thus $\mathcal E_{r_+}(\kappa)\succeq_B\mathcal E_{r_-}(\kappa)$ pointwise under the pointwise
form of (S14), or modulo a prior-null set under its almost-sure form. No positive-mass
condition on either cell is needed. The endpoint values of $\kappa$ refer to the primitive noise
law in (S12).
\end{theorem}
\csname unlabelledfalse\endcsname

\begin{proof}
First, the cells are nested. At the threshold margin, a path that crosses $\tau$ crosses
$\tau'<\tau$ no later. At the clock margin, $c+T\leq H$ implies $c+T'\leq H$. Conditions (S4),
(S5), (S7), and (S11) therefore give
\begin{equation}
\label{eq:blackwell-nesting}
\mathcal C_{F,r_-}\subseteq\mathcal C_{F,r_+},
\qquad
\mathcal C_{P,r_+}\subseteq\mathcal C_{P,r_-}.
\end{equation}

On $\mathcal C_{F,r_+}$, condition (S14) gives
$B^F_{r_+}+Q^F_{r_+}=g(s)$ with $g$ strictly increasing. Its inverse on its image is measurable,
and the Lusin--Suslin theorem makes the image Borel, so $\mathsf S_{F,r_+}$ recovers $s$, then
$j(s)$, and hence $\theta$. Denote this decoder by
$\iota_+$.

Define $K_H$ on the three branches of the tighter output. If $y=(P,\mathfrak h)$, emit
$(P,\mathfrak h)$. If $y=(F,\mathsf s_F)$, recover
$\widehat\theta=\iota_+(\mathsf s_F)$. When $D_{r_-}^H(\widehat\theta)=1$, emit the deterministic
looser tuple $(F,\mathsf S_{F,r_-}(\widehat\theta))$. When
$D_{r_-}^H(\widehat\theta)=0$, draw $z'_{0:H}\sim Q_\kappa$ and emit
\[
\left(P,
\bigl(m_{r_-,0}(\widehat\theta)+z'_0,\ldots,m_{r_-,H}(\widehat\theta)+z'_H;
0,\ldots,0\bigr)\right).
\]
Here $m_{r_-,d}$ denotes the looser rule's mark path.
Define the kernel arbitrarily on flagged tuples outside the image of the tighter flagged channel.
The decoder, disclosure test, tuples, and marks are measurable by (S1) to (S7) and (S14), and the
last branch uses the fixed law in (S12), so this is a Markov kernel that reads only $y$.

Fix a type. If it is pooled under the tighter rule, nesting makes it pooled under the looser rule.
The first branch is exact because (S7), (S11), and (S12) give the same mark path, noise draw, and
all-zero flag coordinates under both rules. If the type is flagged under both, the second branch
emits its looser tuple. If it is newly flagged, the third branch redraws the common noise channel
and adds it to the looser mark path, which is exactly its looser pooled law. These cases prove
\eqref{eq:blackwell-kernel}. If prices are recorded, the kernel discards the input prices and
appends the pinned prices belonging to its emitted output.

Finally, $a$ is a function of $\theta$, the conditional law of $v$ given $\theta$ is
rule-invariant, and the output is conditionally independent of $v$ given $\theta$ by (S7) and
(S12). Integrating \eqref{eq:blackwell-kernel} against the common conditional law of $\theta$
given $(v,a)$ proves the same order about $(v,a)$.
\end{proof}

\csname unlabelledtrue\endcsname
\begin{corollary}[Posterior engagement risk]
\label{cor:blackwell-posterior}
Assume:
\begin{enumerate}[label=(P\arabic*),leftmargin=3.0em,itemsep=3pt]
\item \emph{Standing conditions.} Conditions (S1) to (S14) hold.
\item \emph{Rule comparison.} The two rules satisfy (B2) at a common $\kappa$, so
  Theorem~\ref{thm:blackwell} applies.
\end{enumerate}
Then the expected posterior risk of engagement, measured by posterior variance, falls weakly under
the tighter rule:
\begin{equation}
\label{eq:blackwell-posterior}
\Eop\!\left[\operatorname{Var}(a\mid Y_{r_+})\right]
\leq
\Eop\!\left[\operatorname{Var}(a\mid Y_{r_-})\right].
\end{equation}
The same inequality holds with $a$ replaced by any square-integrable function of $\theta$.
\end{corollary}
\csname unlabelledfalse\endcsname

\begin{proof}
Use $K_H$ to couple the outputs so that $\theta$, $Y_{r_+}$, and $Y_{r_-}$ form a Markov chain in
that order. Put $M_+=\Eop[a\mid Y_{r_+}]$ and $M_-=\Eop[a\mid Y_{r_-}]$. Then
$M_-=\Eop[M_+\mid Y_{r_-}]$. Conditional variance decomposition gives
\[
\Eop\!\left[\operatorname{Var}(a\mid Y_{r_-})\right]
-\Eop\!\left[\operatorname{Var}(a\mid Y_{r_+})\right]
=\Eop\!\left[\operatorname{Var}(M_+\mid Y_{r_-})\right]\geq0.
\]
The same argument applies to every square-integrable function of the type.
\end{proof}

\csname unlabelledtrue\endcsname
\begin{corollary}[The two-parameter Blackwell chain]
\label{cor:blackwell-chain}
Assume:
\begin{enumerate}[label=(K\arabic*),leftmargin=3.0em,itemsep=3pt]
\item \emph{Standing conditions.} Conditions (S1) to (S14) hold at both rules.
\item \emph{Rule comparison.} The pair $r_+,r_-$ satisfies (B2).
\item \emph{Liquidity comparison.} The active order is two noise lumps and
  $0<\kappa_L<\kappa_H<1$, so Lemma~\ref{lem:g2} applies on each pooled cell.
\end{enumerate}
Then
\begin{equation}
\label{eq:blackwell-chain}
\mathcal E_{r_+}(\kappa_L)
\succeq_B\mathcal E_{r_+}(\kappa_H)
\succeq_B\mathcal E_{r_-}(\kappa_H),
\end{equation}
and also
\[
\mathcal E_{r_+}(\kappa_L)
\succeq_B\mathcal E_{r_-}(\kappa_L)
\succeq_B\mathcal E_{r_-}(\kappa_H).
\]
In particular, a tighter rule at lower liquidity dominates a looser rule at higher liquidity.
\end{corollary}
\csname unlabelledfalse\endcsname

\begin{proof}
Fix one rule. Its cell event and flagged tuple do not move with $\kappa$ under (S7), (S10),
(S11), and (S12). Apply the kernel of Lemma~\ref{lem:g2} on the pooled branch and leave the flagged
tuple unchanged. This garbles the tagged experiment at $\kappa_L$ into the one at
$\kappa_H$. If prices are appended, the pooled price is recomputed from the emitted history and
the flagged price is recomputed from the emitted tuple. Theorem~\ref{thm:blackwell}
gives the rule order at either fixed liquidity. Transitivity gives both displayed chains.
\end{proof}

\csname unlabelledtrue\endcsname
\begin{corollary}[When the Blackwell order is strict]
\label{cor:blackwell-strict}
Assume:
\begin{enumerate}[label=(R\arabic*),leftmargin=3.0em,itemsep=3pt]
\item \emph{Standing conditions.} Conditions (S1) to (S14) and the comparison in (B2) hold at a
  common $\kappa\in[0,1]$.
\item \emph{Positive new flagged mass.} The set
  $N=\mathcal C_{F,r_+}\setminus\mathcal C_{F,r_-}$ has positive mass. By the nesting in
  Theorem~\ref{thm:blackwell}, $\Prb(N)=\Omega_{r_+}-\Omega_{r_-}$.
\item \emph{A nonatomic signal on the cut.} The conditional law of $s$ given $N$ is nonatomic.
\item \emph{A finite pooled range.} The looser pooled history takes values in a finite set.
\end{enumerate}
Then the order about $\theta$ in Theorem~\ref{thm:blackwell} is strict: the looser experiment cannot
garble the tighter one. \textsc{Numerical.} On the calibration's threshold ladder, these conditions
hold for each positive-mass cut at $T=5$, so every such threshold comparison is strict. The
threshold cuts at $T=10$ are null, and the corresponding rule experiments are
Blackwell-equivalent. These calibration facts come from
\texttt{numerical\_v4/checks/t2\_threshold\_revelation\_check.json}.
\end{corollary}
\csname unlabelledfalse\endcsname

\begin{proof}
Theorem~\ref{thm:blackwell} gives the weak order. Suppose a reverse kernel existed. By (S7) and
(S13), for every type in $N$, the tighter output is the point mass at its flagged tuple, while the
looser output lies in
the finite pooled range. A mixture of rows of the reverse kernel can equal that point mass only if
every row receiving positive weight is the same point mass. The tighter flagged tuple would
therefore take only finitely many values on $N$ outside the type-null exceptional set on which the
reverse-kernel identity need not hold. Deleting that set preserves $\Prb(N)>0$ and the nonatomic
conditional signal law. Condition (S14) makes that tuple
injective in $s$, while (R3) makes its conditional distribution on $N$ nonatomic. This is a
contradiction, so no reverse kernel exists. The final calibration sentences apply this criterion to
the positive cuts at $T=5$. At the reported null cuts, $T=10=H$ makes every flagged type cross at
date zero, and nesting plus nullity makes the flagged sets coincide up to a prior-null set. On that
set $B^F=b^*(s)$ and $Q^F=0$ under both thresholds, so identity kernels on the flagged and pooled
branches give the stated equivalence.
\end{proof}

\begin{remark}[What the order does not sign]
A Blackwell order at fixed liquidity does not sign the liquidity sensitivity of the engagement
premium $\Delta_m\Eop[h(\mathcal I_H)]$, where $h=\pi p$. It compares level experiments at one
$\kappa$, while changing the rule also changes the type composition and weight of the pooled cell
whose derivative defines that sensitivity.
\end{remark}
